% ─────────────────────────────────────────────────────────────────────────────
% chapters/introduction.tex — General Introduction (unnumbered)
% ─────────────────────────────────────────────────────────────────────────────

\chapter*{General Introduction}
\addcontentsline{toc}{chapter}{General Introduction}
\markboth{General Introduction}{General Introduction}

\section*{General Context}

Algeria's mobile telecommunications market is characterised by a triopoly consisting of three licensed operators: Mobilis, Djezzy, and Ooredoo. Together, these operators serve more than 47 million subscribers, representing a mobile penetration rate that ranks among the highest on the African continent \parencite{ARPCE2024}. Each operator regularly releases prepaid, postpaid, and data-only plans — varying in price, data volume, voice minutes, SMS allowances, validity periods, and network generation (4G or 5G). The resulting catalogue of offers grows continuously and changes frequently, making it increasingly difficult for consumers to identify the plan that best matches their individual needs and budget.

Despite the richness of the available supply, Algerian consumers currently lack access to a neutral, centralised, and up-to-date comparison tool. Operator websites present only their own products, and no independent aggregation platform exists at the national level. This information asymmetry forces users to visit multiple websites, manually record and compare figures, and rely on incomplete or outdated information when making purchasing decisions.

\section*{Problem Statement}

The central problem addressed by this work can be stated as follows: how can a neutral, real-time web platform be designed and implemented to transparently aggregate all Algerian mobile offers, provide advanced filtering and side-by-side comparison tools, and deliver personalised recommendations tailored to each user's budget and usage profile?

Answering this question requires solving three interdependent sub-problems. First, offer data must be collected automatically and reliably, given the absence of public APIs from the operators. Second, a scoring model must translate a user's declared preferences into a ranked list of suitable plans. Third, all of these capabilities must be packaged in a responsive, accessible, and intuitive web interface suitable for a broad audience.

\section*{Objectives and Scope}

This project, carried out within the framework of the Licence en Métiers de l'Informatique at IFAG Alger, pursues five clearly defined objectives. The first objective is to develop a full-stack web application that enables users to browse, filter, and compare mobile plans from all three national operators. The second objective is to build an automated data-collection module capable of scraping offer data from operator websites and maintaining a normalised, up-to-date database. The third objective is to implement a recommendation engine that scores available plans against a user-defined profile and returns the three best-matched options. The fourth objective is to integrate data visualisation tools, including a best-value highlighting system on the comparison table and statistics cards on the admin dashboard, to support informed decision-making. The fifth objective is to implement a real-time notification system that alerts registered users to relevant changes in the offer catalogue.

The scope of this work covers the three Algerian mobile operators (Mobilis, Djezzy, Ooredoo) and the plan types currently offered: prepaid, postpaid, and data-only. Fixed-line internet and bundled services fall outside the scope of this version and are identified as future extensions.

\section*{Methodology}

The project follows an iterative development methodology. Requirements were gathered through a systematic analysis of existing comparison platforms and a direct examination of the three operator websites, identifying both implicit requirements from observed best practices and explicit gaps left unaddressed by the current market. The technical architecture was designed around the \textit{Next.js~16} full-stack framework, which allows both the front-end interface and the back-end API routes to be developed within a single codebase. Data persistence is managed through \textit{Prisma~ORM} backed by a \textit{SQLite} database, chosen for its simplicity and zero-configuration deployment footprint. Automated offer collection is handled by a custom scraping module that runs on a scheduled basis. The recommendation engine applies a weighted multi-criteria scoring algorithm derived from content-based filtering principles.

\section*{Thesis Structure}

The remainder of this mémoire is organised into four chapters. Chapter~1 surveys the state of the art, covering the Algerian telecom market, existing comparison platforms, web data aggregation techniques, recommendation systems, and the modern web stack selected for this project. Chapter~2 presents the requirements analysis and system design, including use case diagrams, the data model, the system architecture, and technology justifications. Chapter~3 describes the full implementation of the platform, detailing the scraping module, the REST API, the recommendation engine, the user interface, and the security mechanisms. Chapter~4 evaluates the platform against the five project objectives, validates the recommendation engine analytically on representative test profiles, documents the identified limitations, and summarises the documentation produced. The work closes with a general conclusion that summarises contributions, acknowledges limitations, and outlines perspectives for future development.

\section*{Note on the Use of AI Writing Assistance}

In accordance with \S{}2.18 of the IFAG Charte de Rédaction du Mémoire, the author declares that the AI language model \textit{Claude} (Anthropic) was used during the preparation of this mémoire. Its assistance was limited to proofreading and reformulating sentences for clarity, verifying technical consistency between the written descriptions and the implemented source code, and correcting factual inaccuracies in diagram labels and table values. All intellectual content, analytical work, design decisions, and implementation are entirely the author's own. The author assumes full responsibility for the accuracy and integrity of this document.
